\documentclass[12pt]{article}
\usepackage[utf8]{inputenc}
\usepackage{amsmath}
\usepackage{amssymb}
\usepackage{amsfonts}
\usepackage{geometry}
\usepackage{enumitem}
\usepackage{chemfig}
\usepackage{multicol}
\geometry{a4paper, margin=1in}

% Adjust chemfig globally for better fit in columns
\setchemfig{atom sep=1.5em, bond offset=1pt, nodesep=2pt}

\begin{document}

\begin{center}
    \subsection*{Alkanes \& Cycloalkanes}
\end{center}

\begin{enumerate}[label=\arabic*.]
    \item What is the general formula for an alkane?
    \vspace{1.5cm}
    \item What is the general formula of a saturated cycloalkane?
    \vspace{1.5cm}
    \item Why are alkanes saturated?
    \vspace{1.5cm}
    \item Are all cycloalkanes saturated? Explain.
    \vspace{1.5cm}
    \item How many hydrogen atoms are in the chemical formula of an alkane that has 3 methyl branches on a 6-carbon chain?
    \vspace{1.5cm}
    \item How many hydrogen atoms are in the chemical formula that has 2 methyl branches on a 5-carbon chain cyclic main chain?
    \vspace{1.5cm}
    \item Why are the MP/BP of alkanes \& cycloalkanes typically low?
    \vspace{1.5cm}
    \item Why does the MP/BP of alkanes \& cycloalkanes increase with mass?
    \vspace{1.5cm}
    \item For structural isomers, why do branched alkanes typically have slightly lower MP/BP than non-branched alkanes?
    \vspace{1.5cm}
    \item Why are alkanes \& cycloalkanes soluble in oil but not in water?
    \vspace{1.5cm}
    \item For structural isomers, why do straight chain alkanes have slightly higher densities than branched alkanes?
    \vspace{1.5cm}
\end{enumerate}

\begin{center}
    \subsection*{Alkynes}
\end{center}

\begin{enumerate}[label=\arabic*.]
    \item What do all alkynes have in common?
    \vspace{1.5cm}
    \item Are alkynes saturated? Explain.
    \vspace{1.5cm}
    \item How many hydrogens are in a structure that has 3 triple bonds in a 8-carbon chain?
    \vspace{1.5cm}
    \item Why are alkynes linear?
    \vspace{1.5cm}
    \item Why do cycloalkynes not really exist?
    \vspace{1.5cm}
\end{enumerate}

\begin{center}
    \subsection*{Geometric Isomers \& Alkenes}
\end{center}

\begin{enumerate}[label=\arabic*.]
    \item What does it mean for a molecule to have a geometric isomer?
    \vspace{1.5cm}
    \item Identify whether each of the following structures is a geometric isomer. If yes, draw the other geometric isomer.
    \begin{multicols}{2}
    \begin{enumerate}[label=\alph*.]
        \item \chemfig{-[:30]=[:-30]-[:30]}
        \item \chemfig{-[:30]-[:-30](=[:60])}
        \item \chemfig{-[:30]=^[:90]-[:150]}
        \item \chemfig{*5(---(-)-)}
        \item \chemfig{-[:30]-[:-30]=[:30]-[:-30]}
        \item \chemfig{-[:30](-[:90])=[:-30]-[:30]}
        \item \chemfig{-[:30](-[:90])=[:-30]}
        \item \chemfig{-[:30]=[:-30]-[:30]-[:-30]}
    \end{enumerate}
    \end{multicols}
    \vspace{2cm}
    \item What do all alkenes have in common?
    \vspace{1.5cm}
    \item Are alkenes saturated? Explain.
    \vspace{1.5cm}
    \item How many hydrogens are in a chemical formula with 2 double bonds on a 5-carbon chain?
    \vspace{1.5cm}
    \item How many hydrogens are in a chemical formula with 2 cyclics, 1 triple bond \& 3 double bonds on a 20-carbon chain?
    \vspace{1.5cm}
\end{enumerate}

\newpage

\begin{center}
    \subsection*{General Physical Properties of Aliphatics}
\end{center}

\begin{enumerate}[label=\arabic*.]
    \item What does it mean for two molecules to be structural isomers?
    \vspace{1.5cm}
    \item What happens to two molecules during a change of state?
    \vspace{1.5cm}
    \item There are three structural isomers of $C_5H_{12}$. Their boiling points vary, with values of $36.1^{\circ}C$, $27.8^{\circ}C$ \& $9.5^{\circ}C$. Draw each isomer, indicate its expected boiling point \& explain its boiling point relative to the other isomers.
    \vspace{4cm}
    \item How would the densities of these three isomers of $C_5H_{12}$ vary? Explain your answer.
    \vspace{2cm}
    \item All hydrocarbons undergo complete combustion reactions to produce $CO_2(g)$, $H_2O(l)$ \& energy. Would the number of $CO_2$ \& $H_2O$ molecules be expected to be the same from burning one gram of each isomer of $C_5H_{12}$? Explain your answer.
    \vspace{2cm}
    \item Would the amount of heat produced be the same when one gram of each isomer of $C_5H_{12}$ undergoes combustion? Explain.
    \vspace{2cm}
    \item Small aliphatic compounds tend to be found in the gas state at room temperature. Explain.
    \vspace{2cm}
    \item When comparing 2-methylpentane with hexane:
    \begin{enumerate}[label=\alph*.]
        \item Which should have the higher boiling point \& why? 
        \item Which should have the higher density \& why?
    \end{enumerate}
    \begin{center}
        \begin{tabular}{cc}
            \chemfig{-[:30](-[:90])-[:-30]-[:30]-[:-30]} & \chemfig{-[:30]-[:-30]-[:30]-[:-30]-[:30]} \\
            2-methylpentane & hexane
        \end{tabular}
    \end{center}
\end{enumerate}

\begin{center}
    \subsection*{Aliphatic Compound Review}
\end{center}
Identify each functional group in each molecule. (Alkane is the default if no other functional group is present.)

\begin{multicols}{3}
\begin{enumerate}[label=\arabic*.]
    \item \chemfig{-[:30]-[:-30]-[:30]-[:-30]}
    \item \chemfig{-[:30]=[:-30]-[:30]=[:-30]}
    \item \chemfig{*4(----)}
    \item \chemfig{-[:30]-[:-30]~[:30]-[:-30]}
    \item \chemfig{-[:30](-[:90])-[:-30](-[:270])-[:30]}
    \item \chemfig{*6(--(-)--(-)-)}
    \item \chemfig{*6(-=-=-=)}
    \item \chemfig{-[:30]-[:-30]=[:30]-[:-30]}
    \item \chemfig{-[:30]-[:-30]-[:30]=[:-30]}
    \item \chemfig{*6(------)}
    \item \chemfig{-[:30]-[:-30]-[:30]-[:-30]=[:30]}
    \item \chemfig{*6(--=---)}
    \item \chemfig{-[:30]-[:-30](-[:270])=[:30]-[:-30]}
    \item \chemfig{-[:30]=[:-30](-[:270])-[:30]=[:-30]}
    \item \chemfig{-[:30]-[:-30]=[:30]-[:-30]-[:30]}
    \item \chemfig{*5(--=-(-)-)}
    \item \chemfig{*6(------)}
    \item \chemfig{-[:30](-[:90])(-[:150])-[:-30](-[:270])-[:30]}
    \item \chemfig{*6(---(-)-(-)-)}
    \item \chemfig{-[:30]-[:-30](-[:270])-[:30](-[:90])-[:-30]}
\end{enumerate}
\end{multicols}

Identify each functional group in each molecule.
\begin{multicols}{2}
\begin{enumerate}[label=\alph*.]
    \item 2,3-dimethylpentane
    \item 2-methylbutane
    \item 5-ethyl-1,3-cyclopentadiene
    \item 3,4-diethyl-2,5-dimethyloctane
    \item 2,3,4-trimethyl-trans-3-hexene
    \item 3,4-diethyl-1-trans-5-octadiene
    \item 4-ethyl-3-methyl-cis-2-heptene
    \item 1,3-dimethylcyclopentane
    \item 5-butyl-2,4,7-trimethyldecane
    \item 3,5-diethyl-1-octyne
    \item 4,5,6-trimethylnonane
    \item 2,3,4,4-tetramethyl-2-hexene
    \item 1,3,5-cyclohexatriene
    \item 3-methyl-1,2-butadiene
    \item 3-methyl-trans-2-hexene
    \item 2-methyl-1-propene
    \item 3-ethyl-1-pentyne
    \item cis-2-pentene
    \item 1,3-cyclooctadiene
    \item 3-methyl-cis-2-pentene
    \item 3,4-dimethyl-1-cyclooctene
    \item 1-trans-3,5-hexatriene
    \item 4-ethyl-3-propyl-1-hexyne
    \item 1,2,3,4-tetramethylcyclobutane
    \item cis-2-trans-4-hexadiene
\end{enumerate}
\end{multicols}

\newpage

\begin{center}
    \subsection*{Benzene-Based (Aromatic) Hydrocarbons}
\end{center}

\begin{enumerate}[label=\arabic*.]
    \item What do all aromatic hydrocarbons have in common?
    \vspace{1.5cm}
    \item When comparing benzene to cyclohexane:
    \begin{enumerate}[label=\alph*.]
        \item Which should have the higher boiling point \& why? 
        \item Which should have the higher density \& why?
    \end{enumerate}
    \begin{center}
        \begin{tabular}{cc}
            \chemfig{*6(-=-=-=)} & \chemfig{*6(------)} \\
            benzene & cyclohexane
        \end{tabular}
    \end{center}
    \vspace{1cm}
    \item Why is the MP/BP of benzene typically low?
    \vspace{1.5cm}
    \item Why is benzene soluble in oil but not in water?
    \vspace{1.5cm}
    \item Why do aromatics have a higher density than alkanes with the same number of carbons?
    \vspace{1.5cm}
\end{enumerate}

\begin{center}
    \subsection*{Alkyl Halides}
\end{center}

\begin{enumerate}[label=\arabic*.]
    \item What do all alkyl halides have in common?
    \vspace{1.5cm}
    \item Which of 1-fluoropropane, 1-chloropropane, 1-bromopropane or 1-iodopropane should have the higher boiling point \& why?
    \begin{center}
        \chemfig{-[:30]-[:-30]-[:30]F} \quad \chemfig{-[:30]-[:-30]-[:30]Cl} \quad \chemfig{-[:30]-[:-30]-[:30]Br} \quad \chemfig{-[:30]-[:-30]-[:30]I} \\
        {\small 1-fluoropropane \quad 1-chloropropane \quad 1-bromopropane \quad 1-iodopropane}
    \end{center}
    \vspace{1cm}
    \item Without looking up the boiling point data for 1-chloropentane \& 3-chloropentane, provide a valid argument supporting why:
    \begin{enumerate}[label=\alph*.]
        \item 1-chloropentane should have the higher boiling point 
        \item 3-chloropentane should have the higher boiling point
    \end{enumerate}
    \begin{center}
        \begin{tabular}{cc}
            \chemfig{-[:30]-[:-30]-[:30]-[:-30]Cl} & \chemfig{-[:30]-[:-30](-[:270]Cl)-[:30]-[:-30]} \\
            1-chloropentane & 3-chloropentane
        \end{tabular}
    \end{center}
    \vspace{1cm}
    \item Why do alkyl halides have higher MP/BP than their corresponding aliphatics (\& aromatics)?
    \vspace{1.5cm}
    \item Why are alkyl halides more soluble in water than aliphatics?
    \vspace{1.5cm}
    \item Why does adding an alkyl halide to water cause the solution to freeze at lower temperatures?
    \vspace{1.5cm}
\end{enumerate}

\begin{center}
    \subsection*{Alcohols}
\end{center}

\begin{enumerate}[label=\arabic*.]
    \item What do all alcohols have in common?
    \vspace{1cm}
    \item Distinguish each OH as $1^{\circ}$, $2^{\circ}$ or $3^{\circ}$.
    \begin{multicols}{2}
    \begin{enumerate}[label=\alph*.]
        \item \chemfig{-[:30]-[:-30]-[:30]OH}
        \item \chemfig{*6(---(OH)--(OH)-)}
        \item \chemfig{-[:30](-[:90]OH)-[:-30]}
        \item \chemfig{-[:30]-[:-30](-[:270]OH)-[:30]-[:-30]OH}
        \item \chemfig{*6(--(-)(OH)---)}
        \item \chemfig{-[:30](-[:90]OH)-[:-30](-[:270]OH)-[:30]OH}
        \item \chemfig{-[:30](-[:90]OH)(-[:150])-[:-30]OH}
        \item \chemfig{HO-[:30]-[:-30]-[:30]OH}
    \end{enumerate}
    \end{multicols}
    \vspace{1cm}
    \item (Challenging): Distinguish each OH as $1^{\circ}$, $2^{\circ}$ or $3^{\circ}$.
    \begin{multicols}{2}
    \begin{enumerate}[label=\alph*.]
        \item 2-methyl-2-propanol
        \item 1,4-dicyclobutyl-2,3-butanediol
        \item phenylmethanol
        \item 1,2-ethanediol
        \item 3-cyclopropyl-2-hydroxyhexane
        \item 2-hexanol
        \item 3-hydroxypentane
        \item 2-hydroxyhexane
    \end{enumerate}
    \end{multicols}
    \vspace{1cm}
    \item How do you know which OH is $1^{\circ}$, $2^{\circ}$ or $3^{\circ}$?
    \vspace{1.5cm}
    \item Assuming both have the same number of carbons, would an alcohol with more OH groups be more soluble in water than alcohols with only one OH group?
    \vspace{1.5cm}
\end{enumerate}

\newpage

\begin{enumerate}[label=\arabic*., start=6]
    \item Explain the MP/BP trend: aliphatics $<$ alkyl halides $<$ alcohols.
    \vspace{1.5cm}
    \item Why does exposure increase MP/BP?
    \vspace{1.5cm}
    \item Why are alcohols generally soluble in water?
    \vspace{1.5cm}
    \item Why is a 1-carbon alcohol more soluble than a 10-carbon alcohol in water?
    \vspace{1.5cm}
    \item Why does adding an alcohol to water cause the mixture to boil at a higher temperature?
    \vspace{1.5cm}
\end{enumerate}

\begin{center}
    \subsection*{Ethers}
\end{center}

\begin{enumerate}[label=\arabic*.]
    \item What do all ethers have in common?
    \vspace{1.5cm}
    \item Which of alcohols or ethers should have the higher boiling point \& why?
    \vspace{1.5cm}
    \item Explain the MP/BP trend: aliphatics $<$ ethers $<$ alcohols.
    \vspace{1.5cm}
    \item Why are ethers good solvents of polar, non-polar \& ionic solutes?
    \vspace{1.5cm}
    \item Why does the solubility of a large ether lower than the solubility of a smaller ether?
    \vspace{1.5cm}
\end{enumerate}

\begin{center}
    \subsection*{Aldehydes}
\end{center}

\begin{enumerate}[label=\arabic*.]
    \item What do all aldehydes have in common?
    \vspace{1.5cm}
    \item Is it necessary to write the number “1” for the location of the carbonyl group in aldehydes? Explain.
    \vspace{1.5cm}
    \item Unlike alkanes, explain why the boiling points of aldehydes do not increase drastically as the number of carbons increase.
    \vspace{1.5cm}
    \item Which of $1^{\circ}$ alcohols or aldehydes should have the higher boiling point \& why?
    \vspace{1.5cm}
    \item Explain the MP/BP trend: aldehydes $<$ alcohols.
    \vspace{1.5cm}
    \item Why are aldehydes less water soluble than alcohols?
    \vspace{1.5cm}
    \item Are aldehydes good general solvents of ionic, polar \& non-polar solutes? Explain.
    \vspace{1.5cm}
\end{enumerate}

\begin{center}
    \subsection*{Ketones}
\end{center}

\begin{enumerate}[label=\arabic*.]
    \item What do all ketones have in common?
    \vspace{1.5cm}
    \item Under what circumstances is it possible to have the name “ethanone” as part of a molecule name?
    \vspace{1.5cm}
    \item Is it necessary to include the location number of the carbonyl group in ketones? Explain.
    \vspace{1.5cm}
    \item Explain the MP/BP trend: for large structural isomers, aldehydes $\doteq$ ketones $<$ alcohols.
    \vspace{1.5cm}
    \item Why are ketones less soluble in water than alcohols?
    \vspace{1.5cm}
    \item Are ketones good general solvents of ionic, polar \& non-polar solutes? Explain.
    \vspace{1.5cm}
\end{enumerate}

\begin{center}
    \subsection*{Carboxylic Acids}
\end{center}

\begin{enumerate}[label=\arabic*.]
    \item What do all carboxylic acids have in common?
    \vspace{1.5cm}
    \item Is it necessary to write the location numbers for the location of the carboxyl group(s) in carboxylic acids? Explain.
    \vspace{1.5cm}
    \item Unlike aliphatics, small carboxylic acids tend to be found in the solid state at room temperature. Explain.
    \vspace{1.5cm}
    \item Explain the MP/BP trend: alcohols $<$ carboxylic acids.
    \vspace{1.5cm}
    \item Why does the MP/BP of carboxylic acids increase with the number of carboxyl groups?
    \vspace{1.5cm}
    \item Why are smaller carboxylic acids more soluble than larger carboxylic acids?
    \vspace{1.5cm}
    \item Why do carboxylic acids tend to be solid state than gas state?
    \vspace{1.5cm}
\end{enumerate}

\newpage

\begin{center}
    \subsection*{Esters}
\end{center}

\begin{enumerate}[label=\arabic*.]
    \item What do all esters have in common?
    \vspace{1.5cm}
    \item Would small esters be found in the gas or solid state at room temperature? Explain.
    \vspace{1.5cm}
    \item Would small esters be soluble in water? Explain.
    \vspace{1.5cm}
    \item Explain the MP/BP trend: esters $<$ carboxylic acids.
    \vspace{1.5cm}
    \item Why are most esters generally insoluble in water?
    \vspace{1.5cm}
    \item Why are esters less polar than carboxylic acids?
    \vspace{1.5cm}
\end{enumerate}

\begin{center}
    \subsection*{Amines}
\end{center}

\begin{enumerate}[label=\arabic*.]
    \item What do all amines have in common?
    \vspace{1cm}
    \item Identify each N as $1^{\circ}$, $2^{\circ}$ or $3^{\circ}$:
    \begin{multicols}{3}
    \begin{enumerate}[label=\alph*.]
        \item \chemfig{-[:30]NH-[:-30]}
        \item \chemfig{H_2N-[:30]-[:-30](-[:270]NH_2)-[:30]}
        \item \chemfig{-[:30]N(-[:90])-[:-30]}
        \item \chemfig{H_2N-[:30]=[:-30]-[:30]}
        \item \chemfig{-[:30]-[:-30](-[:270]NH_2)-[:30]}
        \item \chemfig{-[:30]-[:-30]-[:30]NH_2}
        \item \chemfig{*6(-=-(-N(-[:90])(-[:150]))=-=)}
        \item \chemfig{H_2N-[:30]-[:-30]-[:30]}
        \item \chemfig{H_2N-[:30]=[:-30]-[:30]NH_2}
    \end{enumerate}
    \end{multicols}
    \vspace{1cm}
    \item (Challenging) Identify each N as $1^{\circ}$, $2^{\circ}$ or $3^{\circ}$:
    \begin{multicols}{2}
    \begin{enumerate}[label=\alph*.]
        \item N-cyclopentyl-N-methylhexanamine
        \item 1,2-diaminoethane
        \item 2-aminobutane
        \item 2,2-diaminohexane
        \item benzenamine
        \item N-methylbutanamine
    \end{enumerate}
    \end{multicols}
    \vspace{1cm}
    \item Distinguish between a $1^{\circ}$, $2^{\circ}$ \& $3^{\circ}$ amine.
    \vspace{1.5cm}
    \item Assuming they are structural isomers, which of $1^{\circ}$, $2^{\circ}$ or $3^{\circ}$ amines be the most soluble in water? Explain.
    \vspace{1.5cm}
    \item Assuming they are structural isomers, which of $1^{\circ}$, $2^{\circ}$ or $3^{\circ}$ amines be the most soluble in benzene? Explain.
    \vspace{1.5cm}
    \item Which of 1-aminopropane or 1-hydroxypropane should have the higher boiling point \& why?
    \begin{center}
        \begin{tabular}{cc}
            \chemfig{-[:30]-[:-30]-[:30]NH_2} & \chemfig{-[:30]-[:-30]-[:30]OH} \\
            1-aminopropane & 1-hydroxypropane
        \end{tabular}
    \end{center}
    \vspace{1cm}
    \item Explain why amines have higher MP/BP than aliphatics.
    \vspace{1.5cm}
    \item For amine structural isomers, why do $1^{\circ} > 2^{\circ} > 3^{\circ}$ for MP/BP?
    \vspace{1.5cm}
    \item Why are small amines soluble in water?
    \vspace{1.5cm}
\end{enumerate}

\begin{center}
    \subsection*{Amides}
\end{center}

\begin{enumerate}[label=\arabic*.]
    \item What do all amides have in common?
    \vspace{1.5cm}
    \item Explain why amides are rarely gases at room temperature.
    \vspace{1.5cm}
    \item Why are smaller amides more soluble than larger amides?
    \vspace{1.5cm}
    \item Why do amides tend to be solids or liquids at room temperature?
    \vspace{1.5cm}
\end{enumerate}

\newpage

\begin{center}
    \section*{Organic Chemistry Semi-Unit Review (Part 1) - Structures \& Properties}
\end{center}

\noindent The following is a small selection of questions that, if used appropriately, can better help you prepare for the unit test. This is \underline{not} a complete review sheet. Answer the following questions on a separate piece of paper, using correct format.

\subsection*{Multiple Choice:}
\begin{enumerate}[label=\arabic*.]
    \item The functional group of an amide is:
    \begin{multicols}{3}
    \begin{enumerate}[label=\alph*.]
        \item –OH
        \item –C=C
        \item –NR$_2$
        \item –CONR$_2$
        \item –OR
    \end{enumerate}
    \end{multicols}

    \item Carbon atoms can bond with each other to form:
    \begin{enumerate}[label=\alph*.]
        \item single bonds only
        \item double bonds only
        \item single \& double bonds
        \item single, double \& triple bonds
        \item single, double, triple \& quadruple bonds
    \end{enumerate}
\end{enumerate}

\subsection*{Short Answer \& Full Solution:}
\begin{enumerate}[label=\arabic*.]
    \item Identify the functional group(s) for each. (Alkane is the default if no other functional groups exist on the structure).

    \begin{multicols}{2}
    \begin{enumerate}[label=\alph*.]
        \item \chemfig{*6(---(NH_2)---)}
        \item \chemfig{*6(-=(-Br)-(=F)-=-)}
        \item \chemfig{-[:30](-[:90])-[:-30]-[:30](=[:90]O)-[:-30]OH}
        \item \chemfig{-[:30](=[:90]O)-[:-30]-[:30](=[:90]O)-[:-30]}
        \item \chemfig{-[:30]-[:-30]-[:30]O-[:-30]-[:30]}
        \item \chemfig{-[:30]-[:-30](-[:270])-[:30]-[:-30](=[:270]O)}
        \item \chemfig{-[:30]-[:-30]N(-[:270])-[:30](=[:90]O)-[:-30]}
        \item \chemfig{-[:30]-[:-30](=[:270]O)-[:30]-[:-30]-[:30]}
        \item \chemfig{-[:30](=[:90]O)-[:-30]-[:30](=[:90]O)-[:-30]}
        \item \chemfig{-[:30]-[:-30]-[:30]NH-[:-30]-[:30]}
        \item \chemfig{-[:30](=[:90]O)-[:-30]OH}
        \item \chemfig{-[:30]-[:-30]O-[:30]-[:-30]}
        \item \chemfig{-[:30]-[:-30](=[:270]O)-[:30]N(-[:90])-[:-30]}
        \item \chemfig{*6(-=-(-(=[:90]O)-)=-=)}
        \item \chemfig{-[:30](-[:90])-[:-30](-[:270]OH)-[:30]-[:-30]OH}
        \item \chemfig{-[:30]-[:-30]-[:30]-[:-30](=[:270]O)-[:30]}
        \item \chemfig{-[:30]-[:-30]-[:30](=[:90]O)-[:-30]O-*6(-=-=-=)}
        \item \chemfig{-[:30](-[:90]NH_2)-[:-30](-[:270]NH_2)-[:30]NH_2}
        \item \chemfig{-[:30](-[:90])-[:-30](-[:270]OH)-[:30](-[:90]OH)-[:-30]}
        \item \chemfig{-[:30]~[:-30]-[:30]-[:-30]}
        \item \chemfig{*6(---*6(----)-)}
        \item \chemfig{O=[:30]-[:-30]=[:30]O}
        \item \chemfig{-[:30]-[:-30](-[:270])-[:30](=[:90]O)-[:-30]}
        \item \chemfig{Cl-[:30]-[:-30](=[:270]O)-[:30]Cl}
        \item \chemfig{-[:30]N(-[:90])-[:-30]-[:30]-[:-30]}
        \item \chemfig{-[:30](-[:90])-[:-30]-[:30]-[:-30](=[:270]O)}
    \end{enumerate}
    \end{multicols}
    \begin{multicols}{2}
    \begin{enumerate}[label=\alph*\alph*., start=1]
        \item \chemfig{-[:30]-[:-30]O-[:30]-[:-30]-[:30]}
        \item \chemfig{-[:30]-[:-30](-[:270])-[:30](=[:90]O)-[:-30]}
        \item \chemfig{*6(--O-(=O)--)}
        \item \chemfig{-[:30]-[:-30]=[:30]-[:-30]}
        \item \chemfig{-[:30]-[:-30]=[:30]-[:-30]-[:30]Br}
        \item \chemfig{HO-[:30](=[:90]O)-[:-30]-[:30](=[:90]O)-[:-30]OH}
        \item \chemfig{O=[:30]-[:-30]-[:30]OH}
        \item \chemfig{-[:30]-[:-30]O-[:30]-[:-30]}
        \item \chemfig{NH_2-[:30]-[:-30](-[:270]NH_2)-[:30]-[:-30]NH_2}
        \item \chemfig{-[:30]-[:-30](=[:270]O)-[:30]N(-[:90])-[:-30]}
        \item \chemfig{HO-[:30]-[:-30](-[:270]OH)-[:30]OH}
    \end{enumerate}
    \end{multicols}
\end{enumerate}

\end{document}